\begin{table}[t]
\centering
\caption{Effect of Benefit Cap Reduction on Temporary Accommodation}
\label{tab:main}
\begin{tabular}{lcccc}
\toprule
& (1) & (2) & (3) & (4) \\
\midrule
Cap intensity $\times$ Post & 0.312* & 0.224 & -0.016 & 0.023 \\
& (0.177) & (0.168) & (0.230) & (0.207) \\
\addlinespace
Claimant rate & & -0.006 & & \\
& & (0.016) & & \\
\midrule
LA FE & Yes & Yes & Yes & Yes \\
Year FE & Yes & Yes & Yes & \\
Region $\times$ year FE & & & & Yes \\
Sample & Full & Full & Ex. London & Full \\
Observations & 1870 & 1592 & 1639 & 1870 \\
LAs & 278 & 278 & 245 & 278 \\
Within $R^2$ & 0.013 & 0.008 & 0.000 & 0.000 \\
\bottomrule
\end{tabular}
\begin{tablenotes}[flushleft]\footnotesize
\item \textit{Notes:} Dependent variable is households in temporary accommodation per 1,000 population, measured annually at 31 March (Table 784). Cap intensity is capped households per 1,000 population at May 2017. Post equals one for years 2017--2018. Standard errors clustered at the local authority level in parentheses. * $p < 0.10$, ** $p < 0.05$, *** $p < 0.01$.
\end{tablenotes}
\end{table}
